\documentclass[11.5pt]{scrartcl}

% --- style layer (same template as the author's other notes)
\usepackage{xy-format}
\usepackage{xy-theorem}
\usepackage{fancyhdr}
\usepackage{amsmath}
\addbibresource{n1toolkit.bib}
\pagestyle{fancy}
\fancyhf{}
\fancyfoot[C]{\thepage}
\renewcommand{\headrulewidth}{0pt}
\renewcommand{\footrulewidth}{0pt}

\newcommand{\N}{\mathcal{N}}
\newcommand{\tr}{\operatorname{Tr}}
\newcommand{\Nc}{N_c}
\newcommand{\Nf}{N_f}
% TODO: pick a final name. Working name below.
\newcommand{\toolname}{\textsc{N1obs}}

\title{\toolname: a toolkit for 4d $\N=1$ superconformal observables}

\author[1]{Xingyang Yu}
\affil[1]{Department of Physics, Virginia Tech, Blacksburg, VA 24061, USA}

\date{}

\begin{document}
\maketitle

\begin{abstract}
We present \toolname, an open-source toolkit that computes the standard
protected observables of a four-dimensional $\N=1$ superconformal field theory
directly from its Lagrangian data: the gauge group, the matter content, and
the superpotential. For $SU(N)$ quiver gauge theories, with optional $SU(N)$
global flavor symmetry, the toolkit fixes the superconformal $R$-symmetry by
$a$-maximization (returning \emph{exact} algebraic $R$-charges), and from it
the 't~Hooft anomaly coefficients, the central charges $a$ and $c$, the
Hofman--Maldacena and unitarity diagnostics, and the $\N=1$ superconformal
index. Each quantity is validated against textbook results, including SQCD's
$R = 1-\Nc/\Nf$ with matched flavor anomalies, the exact irrational central
charges of del~Pezzo quivers, and the $s$-confinement index identity of
$SU(2)$ with three flavors. The aim is a small, auditable calculator for these
observables: every output is a closed-form or finite-order quantity with a
stated convention and scope.
\end{abstract}

\section{Introduction}

Many of the quantities one wants from a four-dimensional $\N=1$ superconformal
field theory (SCFT) are fixed by symmetry and anomalies, and are therefore
computable, in principle, by hand from the matter content and superpotential.
In practice the bookkeeping is tedious and error-prone: $a$-maximization
\cite{hep-th/0304128} over the abelian flavor symmetries, the trace formulas
for the central charges $a,c$ \cite{hep-th/9708042}, the various 't~Hooft
anomalies, and the single-letter / matrix-integral evaluation of the
superconformal index \cite{hep-th/0510060,0801.4947,1608.02965}.

\toolname{} is a small Python toolkit that automates this for $SU(N)$ quiver
gauge theories. The input is the Lagrangian data only; the output is the set
of protected observables, each as an exact closed form (or, for the index, a
finite-order series) with an explicit convention. The design goal is
\emph{auditability}: a transparent calculator one can check term by term, with
a clearly stated scope, rather than a black box.

\section{Observables}

The input is a quiver: a set of $SU(N_v)$ gauge nodes, optional $SU(\Nf)$
\emph{global} flavor nodes, bifundamental / adjoint / singlet chiral matter
(arrows), and a single-trace superpotential $W$. \toolname{} computes the
following.

\paragraph{Superconformal $R$-symmetry.}
The trial $R$ runs over the abelian flavor symmetries, i.e.\ the solutions of
$R(W)=2$ that are free of the mixed gauge$^2$--$R$ (ABJ) anomaly
$\sum_i T(r_i)(R_i-1)+T(\mathrm{adj})=0$ at every \emph{gauge} node. The
superconformal $R$ maximizes the trial central charge
\begin{equation}
  a(R)=\tfrac{3}{32}\bigl(3\,\tr R^3-\tr R\bigr)
\end{equation}
\cite{hep-th/0304128}. The returned $R$ satisfies $R(W)=2$ and the
gauge-anomaly conditions exactly by construction; the maximizer is solved
exactly for small flavor dimension and, for larger dimension, recognized as an
exact algebraic number from a high-precision numerical solution with an exact
residual check. Theories with irrational superconformal $R$ are thus returned
in closed form (rational, or an element of a real number field).

\paragraph{Central charges and 't~Hooft anomalies.}
From the superconformal $R$,
\begin{equation}
  a=\tfrac{3}{32}\bigl(3\,\tr R^3-\tr R\bigr),\qquad
  c=\tfrac{1}{32}\bigl(9\,\tr R^3-5\,\tr R\bigr),
\end{equation}
together with the matched 't~Hooft anomalies of the global symmetries: $\tr R$
and $\tr R^3$; the non-abelian flavor anomalies $SU(\Nf)^3$ and
$SU(\Nf)^2\,U(1)_R$ (the flavor central charge being
$k_F=-6\,[SU(\Nf)^2 U(1)_R]$ in the convention of \cite{hep-th/9708042}); and
the abelian flavor anomalies $U(1)^3$, $U(1)^2 U(1)_R$, $U(1)$--grav$^2$, and
$SU(\Nf)^2 U(1)$, where the abelian factors span the kernel of
$\{R(W)=0,\ \text{gauge-anomaly-free}\}$.

\paragraph{Necessary SCFT conditions.}
Unitary 4d $\N=1$ SCFTs obey $a,c>0$ and the Hofman--Maldacena bound
$\tfrac12\le a/c\le\tfrac32$ \cite{0803.1467}; the toolkit flags violations. It
also reports the one-loop beta coefficients
$b_0=3T(\mathrm{adj})-\sum_i T(r_i)$ (a diagnostic only: $b_0<0$ at a node is
allowed, e.g.\ for a free-magnetic-phase node) and scans single-trace gauge
invariants against the chiral-primary unitarity bound $R\ge 2/3$. These are
\emph{necessary} conditions: passing them does not prove the existence of an
interacting fixed point.

\paragraph{Superconformal index.}
The $\mathbb{S}^3\times\mathbb{S}^1$ index
\cite{hep-th/0510060,0801.4947,1608.02965} is assembled from the single-letter
chiral and vector indices, exponentiated plethystically, and projected onto
gauge singlets by the $SU(N)$ matrix integral. The toolkit returns the
unrefined ($p=q$) index, the index refined by abelian flavor fugacities, or
the full two-variable index as a finite-order series in $p,q$.

A single entry point returns the observables:
\begin{verbatim}
    obs = scft_observables(theory, flavor_ranks=[Nf, Nf])
    # R-charges, a, c, Tr R, Tr R^3, Hofman-Maldacena,
    # b0 / asymptotic freedom, flavor and abelian anomalies
    I = index_series(theory, order, flavor_ranks=[Nf, Nf])
\end{verbatim}

\section{Validation}

\paragraph{SQCD.}
For $SU(\Nc)$ with $\Nf$ fundamental flavors, $a$-maximization returns
$R_Q=1-\Nc/\Nf$; the cubic flavor anomaly is $SU(\Nf)^3=\pm\Nc$ (matched across
Seiberg duality \cite{hep-th/9411149}), $SU(\Nf)^2 U(1)_R=-\Nc^2/(2\Nf)$ (so
$k_F=3\Nc^2/\Nf$), the baryonic anomalies $\tr B=B^3=0$ and, in the quark
normalization $B(Q)=+1,\ B(\widetilde Q)=-1$, $B^2 U(1)_R=-2\Nc^2$, and
$b_0=3\Nc-\Nf$. The Hofman--Maldacena gate flags theories below the conformal
window, for example $SU(3)$ with $\Nf=4$ ($a/c=87/247<1/2$), where the naive
computation is not a unitary SCFT; $SU(2)$ with $\Nf=3$ instead saturates
$a/c=1/2$, consistent with its free chiral spectrum below.

\paragraph{Toric quivers and exact irrational central charges.}
For the $dP_0=\mathbb{C}^3/\mathbb{Z}_3$ quiver at the $SU(3)^3$ ranks,
$a=99/16$ and $c=51/8$ (the equal-rank family is $a=3(4N^2-3)/16$,
$c=3(2N^2-1)/8$). For $dP_1=Y^{2,1}$ at the $SU(2)^4$ ranks the toolkit
returns the exact irrational values $a=-\tfrac{739}{4}+52\sqrt{13}$ and
$c=-\tfrac{369}{2}+52\sqrt{13}$ ($a\approx 2.7387$), in agreement with the
$a$-maximization / Sasaki--Einstein volume-minimization correspondence
\cite{hep-th/0503183,hep-th/0506232}.

\paragraph{The index and a duality identity.}
The free-chiral index reproduces the elliptic Gamma function. For $SU(2)$ with
three flavors the gauge integral equals the index of fifteen free chirals at
$R=2/3$ (the $C(6,2)$ gauge invariants of the enhanced $SU(6)$ flavor
symmetry), the $s$-confinement identity of \cite{1608.02965,0910.5944}; the
conifold $SU(2)\times SU(2)$ index correctly counts four mesons and six
baryons at the lowest order.

\section{Scope and limitations}

The toolkit covers $SU(N)$ gauge groups with $SU(N)$ global flavor symmetry and
bifundamental / adjoint / singlet matter with single-trace superpotentials;
$SO/Sp$ gauge groups are a natural extension and are not yet implemented. The
index is a finite-order series (so duality checks are to a chosen order), and
the gauge integral is a direct constant-term extraction whose cost grows with
the rank and the order. $a$-maximization assumes no accidental symmetries and
no operators decoupling below the unitarity bound; the necessary SCFT
conditions are flagged but not proven sufficient.

\section{Code availability}

\toolname{} is open source at \url{https://github.com/xingyang-yu/n1obs}
% TODO: repository URL + license.
with a test suite reproducing the validations above.

\printbibliography[heading=bibliography]

\end{document}
